% ============================================
% DISCUSSION & LIMITATIONS
% ============================================
% 👤 PERSON 5: ML & Evaluation Lead
% ============================================

\section{Discussion and Limitations}
\label{sec:discussion}

In this section, we provide a critical interpretation of the experimental results presented in Section~\ref{sec:results}, analyze the underlying reasons for the system's performance, and acknowledge the constraints of our methodology.

\subsection{Key Insights}

\textbf{Behavioral Features Matter:} 
The Random Forest model's 98.5\% accuracy validates our L1/L2 taxonomy. Feature analysis confirms that \textit{Shannon entropy variance} combined with \textit{sequential write bursts} are the most discriminative indicators. This proves that regardless of API obfuscation, the physical transformation of data into high-entropy output remains a reliable, immutable signal of ransomware activity.

\textbf{Stage-Based Detection:} 
Detecting ransomware only at the encryption phase is often too late. By integrating L1 stages like \textit{Reconnaissance} and \textit{Persistence}, ADRDS flags suspicious activity during the "dwell time" before payload execution. This shifts the defense paradigm from reactive recovery to proactive prevention, allowing automated isolation before encryption keys are generated.

\textbf{Model Trade-offs:} 
While Deep Learning is popular, our results favor Random Forest for this domain. It effectively captures non-linear feature interactions, such as directory traversal linked to entropy spikes, without the "black box" opacity of Neural Networks. RF converged faster and proved more resilient to the noise inherent in system API logs.

\subsection{Limitations}

\textbf{Dataset Constraints:} 
Our evaluation combines Cuckoo Sandbox traces with synthetic simulator data. While diverse, this controlled environment lacks the chaotic "background noise" of real-world workstations. Unpredictable user behaviors, such as gaming or heavy compiling, could introduce stochastic variations not present in training, potentially increasing variance in production environments.

\textbf{Evasion Potential:} 
Sophisticated adversaries could attempt to circumvent detection using "low-and-slow" strategies, spacing out encryption to stay below frequency thresholds. Additionally, "mimicry attacks" that inject benign system calls to dilute malicious sequences pose a challenge. Future work must address adversarial robustness against attacks that statistically mimic benign software distributions.

\textbf{Zero-Day Detection:} 
While generalizing better than signatures, our approach is not immune to concept drift. Novel ransomware using radically different kill chains—such as direct raw I/O access or "Living off the Land" tactics—may initially evade detection. These techniques bypass standard API monitoring, requiring model retraining on new behavioral signatures.

\textbf{False Positives:} 
Distinguishing malicious encryption from legitimate high-intensity tasks is difficult. Applications like video renderers and compression tools often mirror ransomware behavior by rapidly writing high-entropy files, occasionally triggering false alerts. Successful deployment requires a dynamic whitelisting mechanism to mitigate these interruptions and differentiate benign heavy workloads from malicious activity.

\subsection{Generalization}

Leave-One-Group-Out cross-validation, testing on unseen families (e.g., Conti), yielded 92.4\% accuracy. While this suggests strong adaptability, the performance drop from 98.5\% (seen data) indicates some overfitting to specific implementation details. Continuous data collection and retraining pipelines are essential to maintain high efficacy against evolving ransomware strains.

\subsection{Deployment Considerations}

\textbf{Computational Overhead:} 
Real-time feature extraction imposes a negligible 3-5\% CPU overhead on standard workstations. However, this footprint may be prohibitive for resource-constrained endpoints like legacy ICS or low-power IoT devices. Deploying in such environments would require optimizing the Python-based agent into a lower-level language like Rust or C++.

\textbf{Integration Challenges:} 
The prototype uses user-land Windows API hooks. A robust production deployment requires kernel-level monitoring (Mini-Filter drivers) to prevent unhooking by advanced malware. Furthermore, porting ADRDS to heterogeneous enterprise networks involves developing parallel event collection modules for Linux (eBPF) and macOS (Endpoint Security Framework), a significant engineering hurdle.

\subsection{Threats to Validity}

\textbf{Internal Validity:} 
Relying on a simulator for training data introduces potential internal validity threats. If the simulator fails to perfectly replicate kernel-level nuances or race conditions of real malware, the model may learn to detect simulation artifacts. We mitigated this with Cuckoo data, but simulator fidelity remains a critical variable.

\textbf{External Validity:} 
Experiments were conducted on isolated virtual machines with minimal software stacks. This does not fully reflect the chaotic diversity of enterprise networks, including antivirus interference and system lag. Consequently, our reported accuracy metrics should be interpreted as an upper bound of performance in ideal conditions rather than field guarantees.

\textbf{Construct Validity:} 
We prioritized Recall to minimize missed attacks. However, in operational security, the cost of a False Negative (breach) vastly outweighs a False Positive. Standard metrics like Accuracy may not fully capture business risk. Future evaluations should employ cost-sensitive learning metrics to better align model objectives with financial impact.
